% !PRINTED 426
\weberpage{426}
%<<CONT
einsetzt, so findet sich
\[
D = 5^5(108a^5c-135a^4b^2+90a^2bc^2 - 320ab^3c+256b^5+c^4),
\]
so dass sich
\weq{16}{5^5\Delta = a^2D}
ergiebt, und man aus (15) f\"ur $\lambda$ den Ausdruck
\[
\lambda = \frac{11a^3b - ac^2 + 2b^2c + \frac{1}{25}a\sqrt{\frac{1}{5}D}}{2(a^4+abc-b^3)}
\]
erh\"alt, in dem die Quadratwurzel beide Vorzeichen haben kann.

Man sieht, dass ausser der Quadratwurzel aus der Discriminante, die rational durch die Wurzeln ausdr\"uckbar ist, noch $\sqrt5$ darin vorkommt.

Hat man $\lambda$ berechnet, so erh\"alt man aus (11) das Verh\"altniss $z_1:\mu^2$ und aus (12) die letzte Unbekannte $z$ \gesperrt{rational} durch $\lambda$ dargestellt.

\begin{center}
\S.\,113.

\gesperrt{Resolventen sechsten Grades.}
\end{center}

Um die Resolventen $6^{\mathrm{ten}}$ Grades der Ikosa\"edergleichung zu finden, gehen wir von einer der in der Ikosa\"edergruppe enthaltenen Di\"edergruppen $D_5$ aus (\S.\,111, 2.). Die Darstellung der Ikosa\"edergruppe [\S.\,58, (20)] giebt uns unmittelbar die Zerlegung in die Nebengruppen; ist
\weq{1}{Q = \Theta^r,\ \Theta^r\psi,\ r=0,1,2,3,4}
die Di\"edergruppe, von der wir ausgehen, so erhalten wir die volle Ikosa\"edergruppe
\weq{2}{P = Q + Q\chi + Q\chi\Theta + Q\chi\Theta^2 + Q\chi\Theta^3 + Q\chi\Theta^4.}
Ist dann $U_\infty$ eine Function, die durch die Substitutionen von $Q$ unge\"andert bleibt und durch
\[
\chi,\ \chi\Theta,\ \chi\Theta^2,\ \chi\Theta^3,\ \chi\Theta^4
\]
in
\[
U_0,\ U_1,\ U_2,\ U_3,\ U_4
\]
\"ubergeht, so sind die $U_\infty, U_0, U_1, U_2, U_3, U_4$ die Wurzeln einer Resolvente $6^{\mathrm{ten}}$ Grades. (Ueber die Bezeichnung $U_\infty$ vgl. \S.\,65.)
